\documentclass[UTF8,a4paper,12pt]{ctexart}
\usepackage{amsmath,amssymb,amsthm,bm,mathtools}
\usepackage{geometry}
\usepackage{enumitem}
\usepackage{xcolor}
\usepackage{hyperref}

\geometry{left=2.6cm,right=2.6cm,top=2.6cm,bottom=2.6cm}
\hypersetup{
  colorlinks=true,
  linkcolor=blue!55!black,
  citecolor=blue!55!black,
  urlcolor=blue!55!black
}
\setlist[itemize]{leftmargin=2em,itemsep=0.25em,topsep=0.25em}
\setlist[enumerate]{leftmargin=2.2em,itemsep=0.25em,topsep=0.25em}
\sloppy

\newtheorem{problem}{Problem}
\newtheorem{task}{Task}
\newtheorem{remark}{Remark}
\newtheorem{lemma}{Lemma}
\newtheorem{definition}{Definition}

\newcommand{\dd}{\,\mathrm{d}}
\newcommand{\sx}{\sigma_x}
\newcommand{\sy}{\sigma_y}
\newcommand{\eps}{\varepsilon}
\newcommand{\Om}{\Omega}
\newcommand{\PML}{\mathrm{PML}}
\newcommand{\R}{\mathbb{R}}
\newcommand{\Ecal}{\mathcal{E}}
\newcommand{\Dcal}{\mathcal{D}}
\newcommand{\Ltwo}[1]{\left\|#1\right\|_{L^2(\Om)}}
\newcommand{\ip}[2]{\left(#1,#2\right)_{L^2(\Om)}}


\title{Notes: Candidate Energy Ansatz for 2D Maxwell--PML}
\author{Working notes, not a proof}
\date{\today}

\begin{document}
\maketitle

\section{Purpose}

This note lists candidate structures to try.  It should not be cited as a proof.  The main source problem is \texttt{energy\_functional\_problem.tex}.

\section{Naive Energy and Its Defect}

For the displayed model, first take the non-fractional case \(\alpha=1\):
\[
  P+\tau\partial_tP=\chi E,\qquad
  \chi=\epsilon_0(\epsilon_s-\epsilon_\infty).
\]
The naive electromagnetic energy is
\[
  \Ecal_{\rm em}
  =\frac12\int_\Om
  \epsilon_0\epsilon_\infty E^2+\mu H_x^2+\mu H_y^2\,\dd x\dd y.
\]
Under zero boundary flux, its derivative contains
\begin{align*}
&-\epsilon_\infty(\sx+\sy)E^2
-E\partial_tP
-\frac{\sx+\sy}{\epsilon_0}EP
+E(P_{z,1}+P_{z,2})\\
&\quad
-\frac{\mu}{\epsilon_0}(\sy-\sx)H_x^2
-\frac{\mu}{\epsilon_0}(\sx-\sy)H_y^2
+H_xQ_x+H_yQ_y.
\end{align*}
The term \(-E\partial_tP\) is favorable when \(\alpha=1\), because
\[
  E\partial_tP
  =\frac{1}{2\chi}\frac{\dd}{\dd t}\|P\|^2
  +\frac{\tau}{\chi}\|\partial_tP\|^2.
\]
Thus a natural first correction is
\[
  \Ecal_{\rm em+P}
  =\Ecal_{\rm em}+\frac{1}{2\chi}\|P\|^2.
\]
The remaining electric PML terms involving \(EP\), \(EP_{z,1}\), \(EP_{z,2}\), and the magnetic sign-indefinite terms still need auxiliary-variable structure.

\section{Magnetic Shift}

Try
\[
  M_x=Q_x-\frac{\mu}{\eps_0}(\sy-\sx)H_x,\qquad
  M_y=Q_y-\frac{\mu}{\eps_0}(\sx-\sy)H_y.
\]
The model then gives
\[
  \partial_tM_x=\frac{\sy-\sx}{\eps_0}\partial_yE-\frac{\sy}{\eps_0}M_x,
\qquad
  \partial_tM_y=-\frac{\sx-\sy}{\eps_0}\partial_xE-\frac{\sx}{\eps_0}M_y.
\]
These identities suggest trying an energy in \(E,H_x,H_y,P,P_{z,1},P_{z,2},M_x,M_y\).  A possible family is
\[
\begin{aligned}
\Ecal_A
&=\Ecal_{\rm em+P}
 +\frac12\int_\Om
 \alpha_1 P_{z,1}^2+\alpha_2 P_{z,2}^2
 +\alpha_xM_x^2+\alpha_yM_y^2\,\dd x\dd y\\
&\quad
 +\int_\Om
 \beta_1EP_{z,1}+\beta_2EP_{z,2}+\beta_3PP_{z,2}\\
&\quad
 +\beta_xH_xM_x+\beta_yH_yM_y\\
&\quad
 +\gamma_xE\,\partial_y^{-1}M_x
 +\gamma_yE\,\partial_x^{-1}M_y
 \,\dd x\dd y,
\end{aligned}
\]
where the last two nonlocal terms are only a brainstorming device.  Prefer local terms if possible.

\section{Finite-Dimensional Lyapunov Test}

For constant \(\sx,\sy\), Fourier transform in space.  For each \(k=(k_x,k_y)\), write
\[
  \frac{\dd}{\dd t}\widehat U=A(k,\sx,\sy)\widehat U.
\]
Then search for \(G=G(k,\sx,\sy)\) with
\[
  G=G^*>0,\qquad GA+A^*G\le0.
\]
If \(G\) can be chosen independent of \(k\), it is a strong hint for a local \(L^2\) energy.  If \(G\) depends on \(k\), translate it into derivative or pseudo-differential energy only with care.

\section{Cole--Cole Memory}

For \(0<\alpha<1\), expect a positive memory term of the schematic form
\[
  \Ecal_{CC}\sim c_P\|P\|^2+
  \tau_0^\alpha\int_0^\infty w(\lambda)\|\phi(\lambda)\|^2\,\dd\lambda,
\]
but the constants \(c_P\) and weight \(w(\lambda)\) must be derived.  The derivation should use
\[
  P+\tau^\alpha\partial_t^\alpha P
  =\epsilon_0(\epsilon_s-\epsilon_\infty)E,
\]
or an equivalent diffusive representation.  This should be treated only after the \(\alpha=1\) case is understood.

\end{document}
